% GENERATED FILE. Edit canonical lessons, then run npm run generate:books.
% canonical-source-hash: fnv1a64:baa5d84c1445d66e
% canonical-lessons: ML-C33-cintikkuka, ML-C33-manassilaakkuka, ML-C33-vaayikkuka, ML-C33-ezhutuka

\chapter{The Mind and the Page}
\label{ch:ml-mind-and-page}
\hlchaptermodality{33}

\begin{chapteropening}
\textbf{By the end of this chapter:} \emph{I can say that I think, understand, read and write in Malayalam, hear from a verb's own ending which of the four Malayalam borrowed, and put the understander in the dative where Malayalam puts him.}

Say \ml{ഞാൻ} \ml{മലയാളം} \ml{എഴുതുന്നു}, run the chapter's four verbs back and pick out the two wearing the Sanskrit ‑\ml{ഇക്കുക} stamp, and say why Malayalam and Tamil write with one word while Kannada and Telugu, having lost its \ml{ഴ}, reach for another.
\end{chapteropening}

\section[cintikkuka]{\ml{ചിന്തിക്കുക} (cintikkuka) — \textquotedblleft{}to think,\textquotedblright{} on a suffix you have used since Chapter 9}

\label{lesson:ML-C33-cintikkuka}



\begin{quote}
Say \emph{enikku malayāḷaṁ aṟiyāṁ}. Knowing \textbf{arrived} at you, so the \textquotedblleft{}I\textquotedblright{} left the subject slot. Thinking you do — and you keep it.
\end{quote}

\subsection*{What to know first}
\begin{quote}
\textbf{\ml{ചിന്തിക്കുക}} — \emph{cintikkuka} — \textbf{to think}
\end{quote}

\begin{quote}
\textbf{\ml{ഞാൻ} \ml{ചിന്തിക്കുന്നു}.} — \emph{ñān cintikkunnu} — \textquotedblleft{}I think.\textquotedblright{}
\end{quote}

Past \textbf{\ml{ചിന്തിച്ചു}} (\emph{cinticcu}), future \emph{cintikkuṁ}. And \emph{cintikkunnu} is the whole present tense — \emph{ñān}, \emph{nī}, \emph{avar}, one word for everybody.

\subsection*{The letters in this word}
\textbf{\ml{ചി}} (\emph{ca} + the \emph{i} sign) + \textbf{\ml{ന്തി}} (the conjunct \textbf{\ml{ന്ത}} \emph{nta}, also with \emph{i}) + \textbf{\ml{ക്കു}} + \textbf{\ml{ക}}. That \textbf{\ml{ക്ക}} is the held-long \emph{kk} of \emph{kṣamikkuka}.

\begin{cousinweb}
\textbf{\ml{ചിന്ത}} (\emph{cinta}) is a \textbf{noun} first — \textquotedblleft{}thought, care\textquotedblright{} — borrowed whole from Sanskrit \textbf{\dv{चिन्ता}} (\emph{cintā}), which carries the same double sense of thinking and worrying.

Native think-words live on beside it: \textbf{\ml{ഓർക്കുക}} (\emph{ōrkkuka}), narrowed to \textquotedblleft{}remember\textquotedblright{}; \textbf{\ml{നിനക്കുക}} (\emph{ninekka}), now literary.
\end{cousinweb}

\begin{grammarlens}[title={Grammar Lens: ‑\ml{ഇക്കുക}, the suffix that naturalises a loan}]
Chapter 9 made \textbf{\ml{ക്ഷമിക്കുക}} (\emph{kṣamikkuka}, \textquotedblleft{}to forgive\textquotedblright{}) from Sanskrit \textbf{\ml{ക്ഷമ}} plus \textbf{‑\ml{ഇക്കുക}}; Chapter 32 made that ending a passport stamp — \textbf{‑\ml{ഉക}} on an inherited verb, \textbf{‑\ml{ഇക്കുക}} on a borrowing. \emph{Cinta} + \emph{‑ikkuka} is the same machine.

That second slot still takes mood as readily as tense: \emph{cintikkāṁ} (\textquotedblleft{}let's think\textquotedblright{}), \textbf{\ml{ചിന്തിക്കണം}} (\emph{cintikkaṇaṁ}, \textquotedblleft{}must think\textquotedblright{}) on the very \textbf{‑\ml{ണം}} that made \emph{kṣamikkaṇaṁ} polite, \emph{cintikkarutŭ} (\textquotedblleft{}must not think\textquotedblright{}).

Each sister has that suffix in her own shape — Kannada \textbf{‑\ka{ಇಸು}}, Telugu \textbf{‑\te{ఇంచు}} — all inherited from one source. What differs is who used it here. \textbf{Tamil} kept native \textbf{\ta{நினை}} (\emph{niṉai}); \textbf{Telugu} and \textbf{Kannada} built theirs from a word meaning \textquotedblleft{}to say,\textquotedblright{} \emph{anukō}, \textquotedblleft{}to say to oneself.\textquotedblright{} Only Malayalam borrowed — and the word it demoted, \emph{ninekka}, is \emph{niṉai} exactly.
\end{grammarlens}

\subsection*{Your turn}
Say these aloud:

\begin{itemize}
\raggedright
  \item \textquotedblleft{}ñān cintikkunnu,\textquotedblright{} then the past, \textquotedblleft{}cinticcu\textquotedblright{}
  \item noun, verb, and its cousin — \textquotedblleft{}cinta … cintikkuka … kṣamikkuka\textquotedblright{}
  \item the mood ladder — \textquotedblleft{}cintikkāṁ … cintikkaṇaṁ … cintikkarutŭ\textquotedblright{}
  \item three strategies — \textquotedblleft{}niṉai\textquotedblright{} kept, \textquotedblleft{}anukō\textquotedblright{} built, \textquotedblleft{}cintikkuka\textquotedblright{} borrowed
  \item knowing, then thinking — \textquotedblleft{}enikku malayāḷaṁ aṟiyāṁ,\textquotedblright{} \textquotedblleft{}ñān cintikkunnu\textquotedblright{}
\end{itemize}

\begin{tcolorbox}[breakable,colback=teal!4,colframe=teal!35!black,title={Before you move on}]
Say \textquotedblleft{}I think,\textquotedblright{} and name the noun inside it. (\emph{Ñān cintikkunnu}; \textbf{\ml{ചിന്ത}}, Sanskrit for \textbf{thought}.) What does \textbf{‑\ml{ഇക്കുക}} tell you, and which earlier verb wore it? (\textbf{Made from a borrowing}; \emph{kṣamikkuka}.) Say \textquotedblleft{}must think.\textquotedblright{} (\emph{Cintikkaṇaṁ} — the \textbf{‑\ml{ണം}} of \emph{kṣamikkaṇaṁ}, a mood in the tense slot.) Which sister kept a native think-word, and which built one from \textquotedblleft{}say\textquotedblright{}? (\textbf{Tamil}, \emph{niṉai}; \textbf{Telugu}, \emph{anukō}.) Why can \emph{aṟiyuka} not take \emph{ñān}? (\textbf{Knowing arrives at you}.)
\end{tcolorbox}
\section[manassilākkuka]{\ml{മനസ്സിലാക്കുക} (manassilākkuka) — \textquotedblleft{}to understand,\textquotedblright{} or: to put it in the mind}

\label{lesson:ML-C33-manassilaakkuka}



\begin{quote}
Say \emph{ñān cintikkunnu}, then \emph{enikku}. One is what you do; the other is the \textquotedblleft{}to me\textquotedblright{} of Chapter 6. This word needs both.
\end{quote}

\subsection*{What to know first}
\begin{quote}
\textbf{\ml{മനസ്സിലാക്കുക}} — \emph{manassilākkuka} — \textbf{to understand}
\end{quote}

\begin{quote}
\textbf{\ml{ഞാൻ} \ml{മനസ്സിലാക്കുന്നു}.} — \emph{ñān manassilākkunnu} — \textquotedblleft{}I understand.\textquotedblright{}
\end{quote}

Past \textbf{\ml{മനസ്സിലാക്കി}} (\emph{manassilākki}). It is long, but not one lump: you can hear the seams.

\subsection*{The letters in this word}
\textbf{\ml{മ}} (\emph{ma}) + \textbf{\ml{ന}} (\emph{na}) + \textbf{\ml{സ്സി}} (a doubled \textbf{\ml{സ്സ}} \emph{ssa} with the \emph{i} sign) + \textbf{\ml{ലാ}} (\emph{la} + long \emph{ā}) + \textbf{\ml{ക്കു}} + \textbf{\ml{ക}}. Hold that \textbf{\ml{സ്സ}} long.

\begin{cousinweb}
Three pieces, and you have met every one:

\begin{itemize}
\raggedright
  \item \textbf{\ml{മനസ്സ്}} (\emph{manassŭ}) — \textquotedblleft{}\textbf{mind},\textquotedblright{} Sanskrit \textbf{\dv{मनस्}} (\emph{manas})
  \item \textbf{‑\ml{ഇൽ}} (\emph{‑il}) — \textquotedblleft{}\textbf{in},\textquotedblright{} the ending of \emph{vīṭṭil}, \textquotedblleft{}at home\textquotedblright{}
  \item \textbf{\ml{ആക്കുക}} (\emph{ākkuka}) — \textquotedblleft{}\textbf{to make, to put}\textquotedblright{}
\end{itemize}

Put in the mind. That is the whole word — transparent to any Malayalam speaker, which is why it never wore down.

Sanskrit \emph{manas} is from the Indo-European root \emph{\textbf{*men‑}}, \textquotedblleft{}to think,\textquotedblright{} and so are English \textbf{mind} and \textbf{mental} and Latin \emph{mēns} — by a different branch of it, so relatives rather than the same word. (\emph{Mind}'s exact Sanskrit match is \textbf{\dv{मति}} \emph{mati}; \emph{manas}'s is Greek \emph{ménos}.)
\end{cousinweb}

\begin{grammarlens}[title={Grammar Lens: two verbs, one root, and who does the work}]
\textbf{\ml{ആക്കുക}} \emph{makes} something so; its twin \textbf{\ml{ആകുക}} (\emph{ākuka}) means it \emph{becomes} so. Swap them and the sentence changes hands:

\begin{quote}
\textbf{\ml{എനിക്ക്} \ml{മനസ്സിലായി}.} — \emph{enikku manassilāyi} — \textquotedblleft{}I understood.\textquotedblright{}
\end{quote}

Literally: \textquotedblleft{}\textbf{to me} — it became in the mind.\textquotedblright{} No \emph{ñān}. The understander wears the dative \textbf{‑\ml{ഇക്ക്}}, as the knower does in \emph{enikku malayāḷaṁ aṟiyāṁ}. Chapter 6 separated what you \textbf{do} from what \textbf{happens to} you; here one root does both, and the ending decides which. \emph{Manassilāyi} is much the commoner in speech.

Neither form takes \textbf{‑\ml{ഇക്കുക}}: both sit on native \emph{ākuka}, so no passport stamp.
\end{grammarlens}

\subsection*{Your turn}
Say these aloud:

\begin{itemize}
\raggedright
  \item the three pieces — \textquotedblleft{}manassŭ … ‑il … ākkuka\textquotedblright{}
  \item \textquotedblleft{}ñān manassilākkunnu,\textquotedblright{} then the past, \textquotedblleft{}manassilākki\textquotedblright{}
  \item the everyday one — \textquotedblleft{}enikku manassilāyi\textquotedblright{}
  \item three understanders — \textquotedblleft{}enikku … ninakku … avanŭ manassilāyi\textquotedblright{}
  \item the two datives together — \textquotedblleft{}enikku malayāḷaṁ aṟiyāṁ,\textquotedblright{} \textquotedblleft{}enikku manassilāyi\textquotedblright{}
  \item with and without a stamp — \textquotedblleft{}cintikkuka … manassilākkuka\textquotedblright{}
\end{itemize}

\begin{tcolorbox}[breakable,colback=teal!4,colframe=teal!35!black,title={Before you move on}]
Name the three pieces of \emph{manassilākkuka}. (\textbf{Mind}, \textbf{in}, \textbf{put}.) Which English words share its root? (\textbf{Mind}, \textbf{mental}.) Say \textquotedblleft{}I understood\textquotedblright{} the way it is usually said, and name the case you wear. (\emph{Enikku manassilāyi}; the \textbf{dative}, as in \emph{enikku malayāḷaṁ aṟiyāṁ}.) Which of today's two verbs carries \textbf{‑\ml{ഇക്കുക}}, and why? (\emph{Cintikkuka} — its noun was borrowed.)
\end{tcolorbox}
\section[vāyikkuka]{\ml{വായിക്കുക} (vāyikkuka) — \textquotedblleft{}to read,\textquotedblright{} and a mouth that turns out not to be there}

\label{lesson:ML-C33-vaayikkuka}



\begin{quote}
\emph{Manassilākkuka} came apart into three pieces you could name. The next word looks as obliging — and the piece you hear is not there.
\end{quote}

\subsection*{What to know first}
\begin{quote}
\textbf{\ml{വായിക്കുക}} — \emph{vāyikkuka} — \textbf{to read}
\end{quote}

\begin{quote}
\textbf{\ml{ഞാൻ} \ml{മലയാളം} \ml{വായിക്കുന്നു}.} — \emph{ñān malayāḷaṁ vāyikkunnu} — \textquotedblleft{}I read Malayalam.\textquotedblright{}
\end{quote}

Past \textbf{\ml{വായിച്ചു}} (\emph{vāyiccu}). The noun beside it is \textbf{\ml{വായന}} (\emph{vāyana}), \textquotedblleft{}reading.\textquotedblright{}

\subsection*{The letters in this word}
\textbf{\ml{വാ}} (\emph{va} + long \emph{ā}) + \textbf{\ml{യി}} (\emph{ya} + the \emph{i} sign) + \textbf{\ml{ക്കു}} + \textbf{\ml{ക}}. Nothing new since \emph{cintikkuka} — the same \textbf{\ml{ക്കുക}} tail.

\begin{cousinweb}
The obvious guess is \textbf{\ml{വായ്}} (\emph{vāy}), \textquotedblleft{}\textbf{mouth}\textquotedblright{} — reading as mouthing the words. A good guess, and wrong.

Gundert marks the noun \textbf{\ml{വായന}} a \emph{tadbhava} of Sanskrit \textbf{\dv{वच्}} (\emph{vac}), \textquotedblleft{}\textbf{to speak},\textquotedblright{} and gives the Tamil verb as \emph{vācikka}. Sanskrit \textbf{\dv{वाचन}} (\emph{vācana}) is \textquotedblleft{}\textbf{reading aloud}.\textquotedblright{} Reading was named for the \textbf{voice}; the mouth-word only sounds like it.

The \emph{c} did not soften in Malayalam. It softened on the road — a \emph{tadbhava} is a word that came the long way, through the middle Indo-Aryan languages, which drop a single consonant between vowels. Malayalam also holds the same Sanskrit word borrowed \textbf{straight}, \textbf{\ml{വാചകം}} (\emph{vācakam}), so it sits here twice; \emph{vāyana} is the travelled one, and the \textbf{‑\ml{ഇക്കുക}} stamp agrees.

Nor did the sisters agree which word to use. Tamil's everyday \textbf{\ta{படி}} (\emph{paṭi}) is Sanskrit too, from \emph{paṭh‑}; Telugu has \textbf{\te{చదువు}} (\emph{caduvu}). Only \textbf{Kannada} kept a native read-verb in ordinary use, \textbf{\ka{ಓದು}} (\emph{ōdu}) — which Malayalam also has, as \textbf{\ml{ഓതുക}} (\emph{ōtuka}), narrowed to reciting scripture.
\end{cousinweb}

\subsection*{Your turn}
Say these aloud:

\begin{itemize}
\raggedright
  \item \textquotedblleft{}ñān malayāḷaṁ vāyikkunnu,\textquotedblright{} then the past, \textquotedblleft{}vāyiccu\textquotedblright{}
  \item the false trail, then the true one — \textquotedblleft{}vāy\textquotedblright{} … \textquotedblleft{}vac, vācana, vāyana\textquotedblright{}
  \item the same word twice — \textquotedblleft{}vācakam\textquotedblright{} straight, \textquotedblleft{}vāyana\textquotedblright{} travelled
  \item the four read-verbs — \textquotedblleft{}vāyikkuka … paṭi … ōdu … caduvu\textquotedblright{}
  \item read what you know — the planet-days, \textquotedblleft{}tiṅkaḷ … veḷḷi … ñāyar\textquotedblright{} — and the year's first month, \textquotedblleft{}Chiṅṅam\textquotedblright{}
  \item when — \textquotedblleft{}ñān rāvile vāyikkunnu,\textquotedblright{} I read in the morning
  \item what you do, what arrives — \textquotedblleft{}ñān vāyikkunnu,\textquotedblright{} \textquotedblleft{}enikku manassilāyi\textquotedblright{}
  \item the three pieces of understanding — \textquotedblleft{}manassŭ … ‑il … ākkuka\textquotedblright{}
\end{itemize}

\begin{tcolorbox}[breakable,colback=teal!4,colframe=teal!35!black,title={Before you move on}]
Say \textquotedblleft{}I read Malayalam.\textquotedblright{} (\emph{Ñān malayāḷaṁ vāyikkunnu}.) Is \emph{vāy}, \textquotedblleft{}mouth,\textquotedblright{} inside it? (\textbf{No} — the trail runs to Sanskrit \emph{vac}, \textquotedblleft{}speak,\textquotedblright{} and \emph{vācana}, \textquotedblleft{}reciting aloud\textquotedblright{}; the \emph{c} softened \textbf{on the road}, and Malayalam keeps \emph{vācakam} unsoftened beside it.) Do the sisters share a read-verb? (\textbf{No} — \emph{paṭi}, \emph{ōdu}, \emph{caduvu}, and only \emph{ōdu} is native.)
\end{tcolorbox}
\section[eḻutuka]{\ml{എഴുതുക} (eḻutuka) — \textquotedblleft{}to write,\textquotedblright{} and the one letter two sisters kept}

\label{lesson:ML-C33-ezhutuka}



\begin{quote}
Reading came from Sanskrit. Writing did not travel at all — and it kept a sound most of the family has lost.
\end{quote}

\subsection*{What to know first}
\begin{quote}
\textbf{\ml{എഴുതുക}} — \emph{eḻutuka} — \textbf{to write}
\end{quote}

\begin{quote}
\textbf{\ml{ഞാൻ} \ml{മലയാളം} \ml{എഴുതുന്നു}.} — \emph{ñān malayāḷaṁ eḻutunnu} — \textquotedblleft{}I write Malayalam.\textquotedblright{}
\end{quote}

Past \textbf{\ml{എഴുതി}} (\emph{eḻuti}). Note the tail: plain \textbf{‑\ml{ഉക}}, no \textbf{\ml{ഇ}} — no passport stamp, because nothing was borrowed.

\subsection*{The letters in this word}
\textbf{\ml{എ}} (the standing vowel \emph{e}) + \textbf{\ml{ഴു}} (\textbf{\ml{ഴ}} \emph{ḻa} with \emph{u}) + \textbf{\ml{തു}} (\emph{tu}) + \textbf{\ml{ക}} (\emph{ka}). The \textbf{\ml{ഴ}} is the letter of \emph{ēḻu}, \textquotedblleft{}seven\textquotedblright{} — curl the tongue back and let the air pass, neither quite an \emph{l} nor quite an \emph{r}. Standard Malayalam keeps it; northern dialects have let it go, as Kannada did.

\begin{cousinweb}
Malayalam's \textbf{\ml{എഴുതുക}} and Tamil's \textbf{\ta{எழுது}} (\emph{eḻutu}) are not cousins. They are the \textbf{same word} — and the Dravidian etymological dictionary finds it in only two other, much smaller languages. Beside \textquotedblleft{}write, paint, draw\textquotedblright{} it records a stranger sense: \textquotedblleft{}\textbf{become indented by pressure}.\textquotedblright{} A letter is a mark pressed into a leaf with the \textbf{\ml{കൈ}} (\emph{kai}), the hand. (You may see it linked to \textbf{\ml{എഴു}} \emph{eḻu}, \textquotedblleft{}to rise\textquotedblright{}; the dictionary keeps them apart, so treat that as unproven.)

The \textbf{\ml{ഴ}} is the sound inside the name \textbf{\ml{തമിഴ്}} itself. Chapter 7 watched it harden going east: \emph{ēḻu}, \emph{ēḷu}, \emph{ēḍu}. Kannada and Telugu lost the letter — and with it this verb. They write with an unrelated root: \textbf{\ka{ಬರೆ}} (\emph{bare}), \textbf{\te{రాయు}} (\emph{rāyu}), from an old word for \textbf{a drawn line}.

Malayalam has that root too — \textbf{\ml{വരയ്ക്കുക}} (\emph{varaykkuka}) — and gave it the other job: \textbf{to draw}. Two roots, two tasks, as with \emph{uṇṭŭ} and \emph{irikkuka}.
\end{cousinweb}

\subsection*{Your turn}
Say these aloud:

\begin{itemize}
\raggedright
  \item \textquotedblleft{}ñān malayāḷaṁ eḻutunnu,\textquotedblright{} then \textquotedblleft{}eḻuti\textquotedblright{}
  \item when — \textquotedblleft{}ñān rāthri eḻutunnu,\textquotedblright{} I write at night
  \item the two sisters' one word — \textquotedblleft{}eḻutuka … eḻutu\textquotedblright{}
  \item the sound hardening east — \textquotedblleft{}ēḻu … ēḷu … ēḍu\textquotedblright{}
  \item write, then draw — \textquotedblleft{}eḻutuka … varaykkuka\textquotedblright{}
  \item the chapter's \textbf{\ml{നാല്}} — \textquotedblleft{}cintikkunnu, manassilākkunnu, vāyikkunnu, eḻutunnu\textquotedblright{}
  \item with Chapter 32's six that makes \textbf{\ml{പത്ത്}} — \textquotedblleft{}uṇṭŭ, pōkunnu, varunnu, tinnunnu, kāṇunnu, aṟiyunnu\textquotedblright{}
  \item the two that carry the stamp — \textquotedblleft{}cintikkuka, vāyikkuka\textquotedblright{}
  \item and the sentence with no subject — \textquotedblleft{}enikku manassilāyi\textquotedblright{}
\end{itemize}

\begin{tcolorbox}[breakable,colback=teal!4,colframe=teal!35!black,title={Before you move on}]
Say \textquotedblleft{}I write Malayalam,\textquotedblright{} and how far Tamil's word is from it. (\emph{Ñān malayāḷaṁ eḻutunnu}; \textbf{not at all}.) Which sound do the standard languages keep, and what did Kannada and Telugu do about writing? (\textbf{\ml{ഴ}} \emph{ḻ}; they lost it and took another root, \emph{bare}, \emph{rāyu}.) Run the chapter's four, and say which two were borrowed. (\emph{Cintikkunnu, manassilākkunnu, vāyikkunnu, eḻutunnu}; \emph{cinta} and \emph{vac}, wearing \textbf{‑\ml{ഇക്കുക}}.) Which refuses you as its subject? (\emph{Enikku manassilāyi}.)
\end{tcolorbox}
